\scanpage{22}
\mspage{3}

\[
\mathrm{Ob}\,X\simeq F=M/\unsure{\mathcal R|R}.
\]

Si $x,y\in\mathrm{Ob}\,X$, on a
\[
\mathrm{Hom}(x,y)\simeq\{u\in M\mid x=yu\},
\]
la composition des flèches \deleted{\unsure{étant}} donnée par la composition de $M$. La catégorie $X$ admet $a=\dot 1$ comme plus grand élément, tout autre objet $x$ étant de la forme
\[
x=\dot u=au.
\]
\inserted{Notons que $Y$ étant pseudo-cofiltrant, il en est de $=$ de $\deleted{\unsure{Y/F}}X=Y_{/F}$ (en fait, pour \unsure{tt} $F\in\mathrm{Ob}\,Y^\wedge$).}

À quelles conditions $X$ admet-il un plus petit élément $e$, lequel sera de la forme
\[
e=a\psi\qquad(\psi\in M\ \text{fixé})\,?
\]

On voit tout de suite que la condition est la suivante
\[
\tag*{\textcircled{A}}\begin{aligned}&\forall u\in M,\ \exists\,u'\in M\quad\text{tel que}\quad uu'\overset{\mathcal R}\sim\psi,\\ &\bigl(\text{i.e }\exists\,p,q\in\mathbb N\ \text{avec}\quad \varphi^p\psi=\varphi^q uu'\bigr).\end{aligned}
\]

Par exemple, il suffit que $\psi$ soit un élément \uline{minimal} de $M$ (i.e. $\forall u,\exists u'$ avec $\psi=uu'$).

% archive sheet number: 21
